\documentclass[11pt]{article}
\usepackage[margin=1in]{geometry}
\usepackage{amsmath,amssymb,amsthm,mathtools}
\usepackage[hidelinks]{hyperref}
\usepackage{booktabs,array,longtable,pdflscape}

\newtheorem{theorem}{Theorem}[section]
\newtheorem{lemma}[theorem]{Lemma}
\newtheorem{proposition}[theorem]{Proposition}
\newtheorem{corollary}[theorem]{Corollary}
\theoremstyle{remark}
\newtheorem{remark}[theorem]{Remark}
\newcommand{\splus}{s^+}
\newcommand{\sminus}{s^-}
\newcommand{\one}{\mathbf 1}
\newcommand{\Arg}{\operatorname{Arg}}

\title{Positive Square Energy of Every Tricyclic Graph}
\author{AI-generated manuscript; no human author is claimed}
\date{31 July 2026}

\begin{document}
\maketitle

\begin{abstract}
Let $G$ be a finite simple connected graph with $n$ vertices and $n+2$
edges.  We prove $\splus(G)\geq n$, where $\splus$ is the sum of the squares
of the positive adjacency eigenvalues.  The proof first splits the cyclic
blocks by cyclomatic ranks $1+1+1$, $2+1$, or $3$.  The first class is the
tricyclic cactus theorem.  The second consists of one theta block and one
cycle block and follows from an exact doubly nonnegative (DNN) sieve and two
induced-territory repairs.  A rank-three block, after suppressing degree-two
paths, has exactly one of four kernels: a four-path theta, a doubled triangle,
$K_4$, or a doubled $4$-cycle.  We give the correlation dual for the DNN
constant, its exact path lemma, and the finite parity certificates and
structural repairs covering those kernels.  Arbitrary bridges and arbitrary
rooted-tree attachments are included.  We assert no strict inequality in
cases closed only at the DNN boundary.
\end{abstract}

\section{Statement and basic tools}

All graphs are finite, simple, and undirected.  If
$\lambda_1,\ldots,\lambda_n$ are the adjacency eigenvalues, put
\[
 \splus(G)=\sum_{\lambda_j>0}\lambda_j^2,
 \qquad \sminus(G)=\sum_{\lambda_j<0}\lambda_j^2,
 \qquad D(G)=\splus(G)-\sminus(G).
\]
For a tricyclic graph, $\operatorname{tr}A^2=2|E|$ gives
\begin{equation}\label{eq:energy}
 \splus(G)+\sminus(G)=2n+4,
 \qquad \splus(G)=n+2+\frac12D(G).
\end{equation}

\begin{theorem}\label{thm:main}
Every finite simple connected graph $G$ with $|V(G)|=n$ and
$|E(G)|=n+2$ satisfies
\[
 \boxed{\splus(G)\geq n}.
\]
\end{theorem}

We use the following pinching consequence of the variational formula for the
positive part of a symmetric matrix.

\begin{lemma}[Induced superadditivity]\label{lem:super}
If $V(G)=V_1\mathbin{\dot\cup}\cdots\mathbin{\dot\cup}V_r$, then
\[
 \splus(G)\geq\sum_{j=1}^r\splus(G[V_j]).
\]
In particular, for $\sigma(H)=\splus(H)-|V(H)|$ one has
$\sigma(G)\geq\sum_j\sigma(G[V_j])$.
\end{lemma}

\begin{proof}
For a real symmetric matrix $A$,
$\operatorname{tr}(A_+^2)=\max_{X\succeq0}
(2\operatorname{tr}AX-\operatorname{tr}X^2)$.  Insert the block diagonal
matrix whose blocks are the positive parts of the principal matrices
$A[V_j]$.  All cross-block entries of $A$ disappear from the trace, giving
the assertion.
\end{proof}

A nonempty connected tree $T$ is bipartite, so its nonzero spectrum is
symmetric and $\splus(T)=|E(T)|=|V(T)|-1$; hence $\sigma(T)=-1$.
We shall also use the following attached-bicyclic credits.

\begin{lemma}[Bicyclic credits]\label{lem:credits}
Let $H$ be connected and bicyclic, allowing rooted trees at every core
vertex.
\begin{enumerate}
 \item If $H$ is bipartite, then $\sigma(H)=1$.
 \item If its cyclic core is either a two-triangle cactus or
 $\Theta(1,2,2)$, then $\sigma(H)>1$.
 \item Every bicyclic cactus has $\sigma(H)\geq0$; every nonunicyclic cactus
 has $\splus(H)>|V(H)|$.
\end{enumerate}
\end{lemma}

\begin{proof}
The first statement follows from spectral symmetry and $|E(H)|=|V(H)|+1$.
For (2), orient every attached tree toward the core.  Its matching recurrence
replaces the activity $t$ at each core vertex by
\begin{equation}\label{eq:message}
 Q_{u\to v}=t+\sum_{w\text{ child of }u}Q_{w\to u}^{-1}>0,
 \qquad a_v=t+\sum_{u\text{ off core at }v}Q_{u\to v}^{-1}\geq t,
\end{equation}
and extracts one common positive factor.  In the normalized Sachs expansion
$i^{-|V|}\det(itI-A)$, every singleton triangle contributes a strictly
negative imaginary term.  The empty term and the even-cycle or disjoint-pair
terms are real.  Thus the determinant lies in the open lower half-plane for
$t>0$.  Its continuous argument is negative, and the signed Coulson identity
\begin{equation}\label{eq:coulson}
 D(H)=-\frac4\pi\int_0^\infty t\,\Arg
 \bigl(i^{-|V(H)|}\det(itI-A(H))\bigr)\,dt
\end{equation}
gives $D(H)>0$.  Since $H$ is bicyclic,
$\sigma(H)=1+D(H)/2>1$.  Statement (3) is the sharp bicyclic-cactus and
nonunicyclic-cactus theorem; its proof uses the same tree elimination, the
block-additive DNN constant below, and the residual cycle packets
\cite{BicyclicCactus,NonunicyclicCactus}.  We cite it only in its displayed
form.
\end{proof}

\section{The DNN dual and the exact path lemma}

For an edge-containing graph $H$, define
\begin{equation}\label{eq:kappa}
 \kappa(H)=\sup_{\substack{M\succeq0,\ M\geq0\\\one^TM\one>0}}
 \frac{4\bigl(\sum_{uv\in E(H)}\sqrt{M_{uv}}\bigr)^2}
 {\one^TM\one}.
\end{equation}

\begin{lemma}[DNN spectral inequality and correlation dual]
\label{lem:dual}
For every graph $H$,
\begin{equation}\label{eq:dnn-bound}
 \sminus(H)\leq\kappa(H),
 \qquad
 \kappa(H)=\min_{\substack{R\succeq0\\R_{vv}=1}}
 \sum_{uv\in E(H)}\frac{2}{1-R_{uv}},
\end{equation}
where boundary values are interpreted by limits.  Moreover $\kappa$ is
additive under one-vertex sums and $\kappa(T)=|E(T)|$ for a tree.
\end{lemma}

\begin{proof}
Write $A=P-B$ for its positive and negative spectral parts.  Then
$M=B\circ B$ is positive semidefinite and entrywise nonnegative.  Also
$\one^TM\one=\operatorname{tr}B^2=\sminus$ and
$\sminus/2=-\sum_{uv\in E}B_{uv}\leq\sum_E\sqrt{M_{uv}}$.
Equation \eqref{eq:kappa} yields $(\sminus)^2\leq\kappa\sminus$.

For weak duality, weighted Cauchy--Schwarz and a correlation matrix $R$ give
\[
 \one^TM\one-\sum_{uv\in E}2(1-R_{uv})M_{uv}
 =\langle R,M\rangle+2\sum_{uv\notin E}(1-R_{uv})M_{uv}\geq0.
\]
Optimizing the weights proves the ``$\leq$'' part of the dual.  Conversely,
normalize $\langle J,M\rangle=1$ and use
$(\sum_e\sqrt{x_e})^2=\min_{a_e>0,\sum a_e=1}\sum_ex_e/a_e$.
Sion minimax and $(\mathcal S_+\cap\mathcal N)^*=\mathcal S_++\mathcal N$
produce a correlation matrix attaining the reverse inequality; an
$\varepsilon$ truncation gives boundary points.  This proves
\eqref{eq:dnn-bound}.

For a one-vertex sum, fold nonedge entries into diagonal terms by adding
$M_{uv}(e_u-e_v)(e_u-e_v)^T$, then split the Gram vector at the cut vertex
orthogonally between the two sides.  Both the edge numerator and denominator
split, and Cauchy--Schwarz proves subadditivity.  Orthogonally gluing
near-optimizers with the optimal relative scale proves the reverse inequality.
Iterating $\kappa(K_2)=1$ gives the tree formula.
\end{proof}

\begin{lemma}[Exact path elimination]\label{lem:path}
Suppose a kernel edge $uv$ is replaced by a path of length $\ell$.  If the
endpoint correlation is $r=R_{uv}$, its exact excess over the baseline
$\ell$ in the dual is
\begin{equation}\label{eq:path}
 f_\ell(r)=\ell\tan^2\!\left(
 \frac{\arccos((-1)^\ell r)}{2\ell}\right).
\end{equation}
For a fixed $r$ and parity, $f_{\ell+2}(r)<f_\ell(r)$ unless the endpoint
angle vanishes.  Different paths may use the same endpoint vectors, so the
minimum of the sum of these path costs over kernel correlation matrices is
exactly $\kappa-L$, where $L$ is the sum of the path lengths.
\end{lemma}

\begin{proof}
Alternately negate the unit Gram vectors along the path.  The transformed
endpoint angle is $\beta=\arccos((-1)^\ell r)$.  The spherical triangle
inequality says the sum of the $\ell$ successive angles is at least $\beta$;
convexity of $\sec^2(x/2)$ makes equal angles optimal.  Equally spaced planar
vectors on the shortest arc attain the bound, proving exactness.  With
$h_s(x)=s\tan^2(x/s)$ and $z=x/s$,
\[
 \frac{d}{ds}h_s(x)=\tan z\,(\tan z-2z\sec^2z)<0
\]
because $\sin z\cos z<z$.  This proves parity monotonicity.
\end{proof}

If a rank-three core has $L$ edges, it has $L-2$ vertices.  Thus
\begin{equation}\label{eq:threshold}
 \kappa\leq L+2\quad\Longrightarrow\quad
 \splus=2L-\sminus\geq L-2=|V|.
\end{equation}
By Lemma~\ref{lem:dual}, the same implication survives arbitrary rooted-tree
attachments: each attached tree edge adds one to both $\kappa$ and the vertex
count.

\section{Cyclic-block classification}

We use the standard block decomposition: a block is a maximal connected
subgraph with no cut vertex, with bridges regarded as two-vertex blocks. Thus
a cycle meeting the rest of the graph in a single vertex is a separate
rank-one block. All degrees used when suppressing a cyclic block are degrees
within that block.

The cyclomatic rank of a connected graph is $|E|-|V|+1$.  Bridge blocks have
rank zero, and ranks add over the nontrivial blocks.

\begin{proposition}[Block alternatives]\label{prop:blocks}
The positive ranks of the cyclic blocks of a connected tricyclic graph are
exactly one of
\[
 1+1+1,\qquad 2+1,\qquad 3.
\]
Rank-one blocks are cycles.  A simple rank-two $2$-connected block is a theta.
A rank-three $2$-connected block, after suppressing each maximal path whose
internal vertices have degree two in the block, has exactly one of the
following loopless multigraph kernels:
\begin{enumerate}
 \item four parallel edges between two branch vertices;
 \item a triangle with two adjacent sides doubled;
 \item $K_4$;
 \item a $4$-cycle with two opposite sides doubled.
\end{enumerate}
Every omitted part is a bridge path or a rooted tree attached at one vertex.
\end{proposition}

\begin{proof}
Block ranks are positive integers summing to three, giving the three
partitions.  A rank-one $2$-connected graph is a cycle.  In rank two,
$\sum_v(d(v)-2)=2$ after suppressing degree two; $2$-connectivity forces two
degree-three vertices joined by three parallel kernel edges.

Suppression cannot create a loop. Such a loop would come from a cycle based at
a branch vertex $u$ whose other vertices all have degree two in the block. If
the block had any part outside that cycle, deleting $u$ would separate the
cycle interior from that part, contrary to $2$-connectivity. If it had no
outside part, the block would itself be a cycle and have rank one.

In rank three the same sum is four.  Hence the branch degrees are either
$(4,4)$, $(4,3,3)$, or $(3,3,3,3)$.  The first gives four parallel edges.
For $(4,3,3)$, looplessness and $2$-connectivity force the degree-four vertex
to have two edges to each degree-three vertex, with one edge between the latter:
 the doubled triangle.  A cubic loopless $2$-connected multigraph on four
vertices is either simple, hence $K_4$, or has parallel edges.  Choose a
parallel pair $uv$.  It cannot be a triple edge, since then $\{u,v\}$ would
be a component, so it is exactly double.  Each of $u,v$ has one remaining
incidence.  If both met the same remaining vertex $x$, the fourth vertex $y$
could meet only $x$; making $d(y)=3$ would already give three $xy$ edges and
force $d(x)\geq5$.  Thus the two incidences meet the two remaining vertices
$x,y$ separately.  Now $x$ and $y$ each have residual degree two, and their
only available mutual edge therefore occurs twice.  Consequently the kernel
has the double edges $uv,xy$ and two single connectors, hence is a $4$-cycle
with opposite sides doubled.  Simplicity of the original graph says at most
one path in each parallel family may have length one.  Reversing suppression
and leaf deletion gives the final assertion.
\end{proof}

\section{The $1+1+1$ and $2+1$ blocks}

\begin{proposition}\label{prop:111}
Every tricyclic graph with block ranks $1+1+1$ satisfies $\splus>n$.
\end{proposition}

\begin{proof}
Such a graph is a cactus of cyclomatic rank three, including all bridge paths
and tree attachments.  Apply Lemma~\ref{lem:credits}(3).
\end{proof}

For a theta with path lengths $a,b,c$, put $L=a+b+c$ and
\begin{equation}\label{eq:Delta}
 \Delta(a,b,c)=\min_{0\leq\theta\leq\pi}
 \left(\sum_{\ell\ \mathrm{even}}\ell\tan^2\frac{\theta}{2\ell}
 +\sum_{\ell\ \mathrm{odd}}\ell\tan^2\frac{\pi-\theta}{2\ell}\right).
\end{equation}
For a cycle put $\epsilon_q=0$ for even $q$ and
$\epsilon_q=q\tan^2(\pi/(2q))$ for odd $q$.

\begin{lemma}[Theta--cycle DNN sieve]\label{lem:21sieve}
For a graph whose cyclic blocks are $\Theta(a,b,c)$ and $C_q$,
\[
 \splus(G)\geq n+2-\Delta(a,b,c)-\epsilon_q.
\]
This is at least $n$ except possibly for
\[
 \Theta(1,2,r)+C_3\quad(r\geq2),
 \qquad \Theta(1,2,2)+C_5.
\]
\end{lemma}

\begin{proof}
Lemma~\ref{lem:path} gives $\kappa(\Theta)=L+\Delta$ and
$\kappa(C_q)=q+\epsilon_q$; one-vertex additivity includes every intervening
bridge and attached tree.  The displayed spectral bound follows from
Lemma~\ref{lem:dual}.  We now give the numerical classification needed for
the sieve.  Put $x=\theta/2$ and
\[
 h_s(y)=s\tan^2(y/s).
\]
For fixed $0<y<\pi/2$ this is strictly decreasing in $s$, since, with
$z=y/s$,
\begin{equation}\label{eq:hs-derivative}
 \frac{d}{ds}h_s(y)
 =\frac{\tan z}{\cos^2z}\bigl(\sin z\cos z-2z\bigr)<0.
\end{equation}
Thus, within a fixed parity, replacing an even length by $2$ or an odd
length by $1$ or $3$ can only increase the objective in
\eqref{eq:Delta}.  Because the graph is simple, at most one of the three
theta paths has length one.

First suppose that no path has length one, and let $e$ be the number of even
paths.  Parity monotonicity bounds the objective from above by
\[
 F_e(x)=2e\tan^2(x/2)+3(3-e)
             \tan^2\bigl((\pi/2-x)/3\bigr).
\]
For $e=0$ and $e=3$, take respectively $x=\pi/2$ and $x=0$, obtaining zero.
For the two mixed cases the explicit tests
\[
 F_1(\pi/4)<1,
 \qquad F_2(\pi/6)<1
\]
follow from $\tan(\pi/8)=\sqrt2-1$,
$\tan(\pi/12)=2-\sqrt3$, and
$\tan(\pi/9)<\tan(\pi/8)$; after substitution the two upper bounds are
$48-4\sqrt2-24\sqrt3$ and $37-6\sqrt2-16\sqrt3$, both less than one by
successive squaring of positive sides.  Hence $\Delta<1$ whenever there is
no length-one path.

Now suppose that there is one such path.  Its contribution is $\cot^2x$.
If the other two paths are odd, the choice $x=\pi/2$ gives $\Delta=0$.
If exactly one of them is even and its length is at least $4$, monotonicity
and the test $x=\pi/3$ give
\begin{equation}\label{eq:one-unit-mixed-test}
 \Delta\leq \frac13+4\tan^2\frac{\pi}{12}
                    +3\tan^2\frac{\pi}{18}<1.
\end{equation}
Indeed $\tan(\pi/18)<\tan(\pi/12)<1/3$, and the resulting inequality follows
from $48\sqrt3>83$, whose square is $6912>6889$.  If both remaining paths
are even and both have length at least $4$, the same test gives
\begin{equation}\label{eq:one-unit-even-test}
 \Delta\leq\frac13+8\tan^2\frac{\pi}{12}<1,
\end{equation}
because $12(2-\sqrt3)^2<1$, equivalently $6889<6912$ after squaring.

It remains precisely to consider sorted triples $(1,2,r)$, $r\geq2$.
The first two path terms satisfy, on writing $u=\tan(x/2)$,
\begin{equation}\label{eq:theta-base-one}
 \cot^2x+2\tan^2(x/2)
 =\frac{1}{4u^2}-\frac12+\frac94u^2\geq1.
\end{equation}
The third term is positive in the interior; at whichever endpoint it can
vanish, the expression in \eqref{eq:theta-base-one} is either divergent or
equal to $2$.  Compactness therefore makes the minimum strictly larger than
$1$.  This proves the exact classification
\begin{equation}\label{eq:Delta-classification}
 \Delta(a,b,c)>1
 \quad\Longleftrightarrow\quad
 (a,b,c)=(1,2,r)\ \text{after sorting},\quad r\geq2.
\end{equation}
For $r=2$, the whole objective is
\[
 \frac{1}{4u^2}-\frac12+\frac{17}{4}u^2,
\]
so minimization at $u^2=1/\sqrt{17}$ gives the exact maximum
\begin{equation}\label{eq:Delta-122}
 \Delta(1,2,2)=\frac{\sqrt{17}-1}{2}.
\end{equation}
For even $r\geq4$, \eqref{eq:hs-derivative} makes the objective pointwise
smaller than for $r=2$.  Moreover, away from $r=2$ one has the uniform strict
bound $\Delta(1,2,r)<4/3$: for odd $r\geq3$, test $x=\pi/3$ and use
$\tan(\pi/18)<1/3$, while for even $r\geq4$ the same test reduces to
$1+4(2-\sqrt3)^2<4/3$, already contained in $6889<6912$.  Since
$4/3<(\sqrt{17}-1)/2$, this also proves that \eqref{eq:Delta-122} is the
global maximum.

Finally, the odd-cycle excess is $h_q(\pi/2)$, so
\eqref{eq:hs-derivative} says that it decreases strictly with odd $q$.  The
first two values and a convenient bound for the rest are
\begin{equation}\label{eq:cycle-excess-cases}
 \epsilon_3=1,
 \qquad \epsilon_5=5-2\sqrt5<\frac23,
 \qquad \epsilon_q\leq\epsilon_7<\frac25\quad(q\geq7\text{ odd}).
\end{equation}
Here the first two identities are exact.  For the last inequality, use
$\sin t<t$, $\cos t>1-t^2/2$, and $\pi<22/7$ at $t=\pi/14$; clearing
denominators gives $7\tan^2t<2/5$.  Even cycles have zero excess.

We need $\Delta+\epsilon_q\leq2$.  For $q=3$,
\eqref{eq:Delta-classification} leaves exactly every $(1,2,r)$.  For $q=5$,
the bound $\Delta<4/3$ disposes of all those triples except $(1,2,2)$, and
\eqref{eq:Delta-122} shows that this remaining triple really exceeds the DNN
threshold: $(\sqrt{17}-1)/2+5-2\sqrt5>2$, since
$\sqrt{17}+7>4\sqrt5$.  For odd $q\geq7$,
$(\sqrt{17}-1)/2<8/5$ and \eqref{eq:cycle-excess-cases} give a strict total
below $2$; even $q$ is immediate.  Thus the exact residuals of this DNN sieve
are the two rows stated.  They are failures of the DNN threshold, not graph
counterexamples.
\end{proof}

\begin{proposition}\label{prop:21}
Every tricyclic graph with block ranks $2+1$ satisfies $\splus(G)>n$.
\end{proposition}

\begin{proof}
Only the rows of Lemma~\ref{lem:21sieve} need repair.  In
$\Theta(1,2,r)+C_3$, choose an internal theta vertex not used by the unique
block-tree route to the external triangle: for $r=2$ choose one of the two
length-two internal vertices; for $r\geq3$ choose one of two internal vertices
of the $r$-path.  Delete that vertex together with its rooted branch.  The
deleted induced territory is a nonempty tree of credit $-1$, while the
connected complement is an attached two-triangle bicyclic cactus of credit
$>1$ by Lemma~\ref{lem:credits}.  Lemma~\ref{lem:super} gives positive total
credit.

For $\Theta(1,2,2)+C_5$, a separating bridge gives two induced territories.
The theta territory has credit $>1$ and the attached pentagon has credit at
least $-(\sqrt5-2)$, by the one-cycle phase comparison
$D\geq-2(\sec(\pi/5)-1)$.  Their sum is positive.  If the blocks share a cut
vertex, keep the theta and that cut vertex in one territory and put the other
four pentagon vertices, with all their rooted branches, in the other.  The
latter is a nonempty tree of credit $-1$ and the former has credit $>1$.
These assignments cover every bridge length, cut location, and rooted branch.
\end{proof}

\section{The four rank-three kernels}

\subsection{Four-path theta}

\begin{proposition}\label{prop:fourpath}
Every simple subdivision of the four-parallel-edge kernel, with arbitrary
rooted trees, satisfies $\splus\geq n$.
\end{proposition}

\begin{proof}
For path lengths $\ell_1,\ldots,\ell_4$, Lemma~\ref{lem:path} reduces the
excess to
\[
 \Phi(\theta)=\sum_{\ell\ \mathrm{even}}\ell\tan^2\frac{\theta}{2\ell}
 +\sum_{\ell\ \mathrm{odd}}\ell\tan^2\frac{\pi-\theta}{2\ell}.
\]
If no length is one, monotonicity reduces even lengths to two and odd lengths
to three.  Writing $x=\theta/2$ and letting $e$ be the number of even paths,
test
$2e\tan^2(x/2)+3(4-e)\tan^2((\pi/2-x)/3)$ at an endpoint for
$e=0,1,3,4$; for $e=2$ its right derivative at $x=0$ is negative.  The
minimum is at most two.

If one length is one (there cannot be two by simplicity), its term is
$\cot^2x$.  For zero, one, or two even paths among the remaining three,
$x=\pi/2$ or $x=\pi/3$ gives a value below two.  With three even paths,
\[
 \Phi(2x)\leq\cot^2x+6\tan^2(x/2).
\]
At $u=\tan(x/2)$ the right side is
$1/(4u^2)-1/2+25u^2/4$, with minimum two at $u^2=1/5$.
Strict path monotonicity improves this unless the lengths are $(1,2,2,2)$,
where direct substitution gives equality.  Hence $\kappa\leq L+2$, and
\eqref{eq:threshold} proves the claim.  Notice that the exceptional core has
spectrum $\{3,-2,-1,0,0\}$ and actually has $\splus=9>5$; we do not infer a
general strict statement from this isolated computation.
\end{proof}

\subsection{Finite kernel certificates}

The remaining arguments repeatedly use the following finite-certificate
principle.  Encode a path parity by $0$ for even and $1$ for odd.  Switching at
a branch vertex toggles all incident bits.  For fixed physical parities,
Lemma~\ref{lem:path} reduces all lengths to the first simple lengths: two for
even, one for odd, except that two parallel odd paths use $1,3$.  A displayed
correlation Gram matrix with path excess at most two then proves every longer
realization of that physical row.  A switching label is used only to organize
rows; it is never used to transport a certificate between physical rows with
different first lengths.

\begin{lemma}[Doubled $4$-cycle certificate]\label{lem:dc4}
Every simple subdivision of the doubled-$4$-cycle kernel satisfies
$\kappa\leq L+2$, except for the structural rows explicitly repaired in the
proof; all rows nevertheless satisfy $\splus\geq n$.
\end{lemma}

\begin{proof}
Write the paths as $a,A:01$, $b:12$, $c,C:23$, $d:30$.
The eight switching invariants and the complete numbers of simple physical
rows are shown in Table~\ref{tab:dc4}.  Exact planar Gram vectors
$u_j=(\cos(\pi k_j/6),\sin(\pi k_j/6))$ give excess at most two in the first
seven classes; all simple representatives are checked separately.  The two
boundary certificate types have contributions, in a suitable order,
$2/3,1/3,0,2/3,1/3,0$, summing to two; all other listed certificates have
strict slack.  This is the finite table proved by direct substitution in
\eqref{eq:path} in \cite{DoubledC4}.

\begin{table}[ht]
\centering\small
\begin{tabular}{cccc}
\toprule
invariant & first simple row & physical rows & disposition\\
\midrule
000&000000&2&DNN\\
001&000001&2&DNN\\
010&000010&4&DNN\\
011&000011&4&DNN\\
100&010000&4&DNN\\
101&010001&4&DNN\\
110&010010&8&DNN, boundary allowed\\
111&010011&8&DNN or deletion\\
\bottomrule
\end{tabular}
\caption{Doubled-$4$-cycle switching census; the counts total all $36$ simple
physical parity rows.}\label{tab:dc4}
\end{table}

For completeness, normalize class $111$ so each doubled side has one even and
one odd path and the connectors have opposite parity.  If a member of a
doubled pair is longer than canonical, the exact planar certificates at
lengths $(4,1,1,2,1,2)$ and $(2,3,1,2,1,2)$ have respective excesses bounded
by
\[
 4\tan^2\frac{3\pi}{32}+3\tan^2\frac\pi8
 +2\tan^2\frac\pi{16}+2\tan^2\frac{3\pi}{16}<1.862,
\]
\[
 2\tan^2\frac\pi{16}+4\tan^2\frac\pi8
 +2\tan^2\frac{3\pi}{16}<1.662.
\]
If both doubled pairs are canonical, open an internal vertex of the even
connector.  The deleted territory is a tree of credit $-1$; the complement is
a two-triangle cactus of credit $>1$.  This exhausts class $111$ and includes
all attachments.
\end{proof}

\begin{lemma}[Doubled-triangle certificate]\label{lem:dtri}
Every simple subdivision of the doubled-triangle kernel, with arbitrary rooted
trees, satisfies $\splus\geq n$.
\end{lemma}

\begin{proof}
Write $a,A:01$, $b,B:02$, and $c:12$.  A doubled pair has first-simple type
$EE=(2,2)$, $EO=(1,2)$, or $OO=(1,3)$.  Table~\ref{tab:dtri} is the complete
$32$-row physical census.  In eleven orbits an exact rational $3\times3$
correlation matrix gives excess at most two.  The entries are obtained from
multiple-angle formulas; for example, with parameter $t$,
\[
 r_{EE}=r_{EO}=\frac{1-6t^2+t^4}{(1+t^2)^2},\quad
 e_{EE}=4t^2,\quad e_{EO}=\left(\frac{1-t^2}{2t}\right)^2+2t^2,
\]
and, for $q=(3t-t^3)/(1-3t^2)$,
$r_{OO}=-(1-q^2)/(1+q^2)$ and $e_{OO}=q^2+3t^2$.
Substitution, nonnegative determinant, and the exact excess fractions are
listed in \cite{DoubledTriangle}; thus no floating optimization is used.

\begin{table}[ht]
\centering\small
\begin{tabular}{ccccc}
\toprule
pair types&rows for even $c$&rows for odd $c$&DNN rows&structural rows\\
\midrule
$EE,EE$&1&1&2&0\\
$EE,EO$&4&4&8&0\\
$EE,OO$&2&2&4&0\\
$EO,EO$&4&4&4&4\\
$EO,OO$&4&4&8&0\\
$OO,OO$&1&1&2&0\\
\midrule
total&16&16&28&4\\
\bottomrule
\end{tabular}
\caption{Doubled-triangle physical parity census.}\label{tab:dtri}
\end{table}

The four structural rows are $EO,EO$ with odd $c$.  If any parallel path is
noncanonical, exact planar certificates reduce to $(4,1,2,1,1)$ or
$(2,3,2,1,1)$ and have excess below $229/120$ or $31/20$.  Otherwise both
parallel pairs are $\{1,2\}$.  If $c\geq3$, open $c$ and leave a two-triangle
cactus; if $c=1$, open the even member of one pair and leave
$\Theta(1,2,2)$.  Lemmas~\ref{lem:super} and \ref{lem:credits} pay the deleted
tree.  These alternatives exhaust every length and attachment.
\end{proof}

\begin{lemma}[$K_4$-subdivision certificate]\label{lem:k4}
Every simple subdivision of $K_4$, with arbitrary rooted trees, satisfies
$\splus\geq n$.
\end{lemma}

\begin{proof}
Switching the six path parities has eight classes, each containing eight
physical rows.  Label the branch vertices $1,2,3,4$, switch so that the three
paths to vertex $4$ are odd, and let $\epsilon_{12},\epsilon_{13},
\epsilon_{23}$ indicate which normalized triangle paths are even.  The word
$000$ is the all-odd class.  For a word of weight one use, up to a permutation
of $1,2,3$, the normalized planar angles
$(2\pi/3,4\pi/3,4\pi/3,0)$; for weight two use
$(2\pi/3,4\pi/3,\pi,0)$; for weight three use $(\pi,\pi,\pi,0)$.
After undoing the switch on the vectors, but not on the physical path lengths,
every canonical physical path has cost $0$, $a=2\tan^2(\pi/12)=14-8\sqrt3$,
or $b=\tan^2(\pi/6)=1/3$.  The complete eight-switch substitution for either
a weight-one or a weight-two class is Table~\ref{tab:k4}; permutations give
all six such classes.  The weight-three template has zero cost in all eight
rows.  This is the exact physical-row table of \cite{K4Seven}, not a transport
of canonical lengths across a switch.

\begin{table}[ht]
\centering\small
\begin{tabular}{cccc}
\toprule
number of rows&number of $a$ terms&number of $b$ terms&total excess\\
\midrule
1&0&5&$5/3$\\
2&2&3&$2a+1$\\
4&3&2&$3a+2/3$\\
1&4&1&$4a+1/3$\\
\bottomrule
\end{tabular}
\caption{Exact eight-switch physical-cost census for each weight-one or
weight-two $K_4$ class.  There are three classes of each weight; the remaining
weight-three class has eight zero-cost rows.  Thus the table and that last
class account for all $6\cdot8+8=56$ non-all-odd physical rows.}
\label{tab:k4}
\end{table}

Exactly, $a<1/6$ because $\sqrt3>83/48$, whose square is the integer
inequality $6912>6889$.  Hence the four totals in Table~\ref{tab:k4} are at
most $5/3<2$.  Fixed-physical-parity path monotonicity proves all longer
realizations, and one-vertex additivity includes all rooted trees.  In
particular, the seven-class ledger proves the strict bound $\kappa<L+2$ in
all $56$ rows.  The normalized templates, all eight switch triples, the
four-type physical substitution table, and these exact inequalities are
given explicitly in \cite{K4Seven}.

Here is the independent all-length cover of the last class
\cite{K4Cover}.  Call an odd path
long when its length is at least three.  With at least three long paths, use
four regular-simplex vectors, $R_{ij}=-1/3$: a unit path contributes $1/2$
and a long path contributes at most
$3\tan^2(\arccos(1/3)/6)<1/6$, so the total is at most two.  With exactly two
long paths, there are two placements.  Opposite paths use branch angles
$(0,\pi/4,\pi,5\pi/4)$ and give
$8\tan^2(\pi/8)=24-16\sqrt2<2$.  Adjacent paths use
$(0,3\pi/8,13\pi/8,\pi)$ and give
\[
 6\tan^2\frac{5\pi}{48}+\tan^2\frac\pi8
 +2\tan^2\frac{3\pi}{16}<\frac{71}{40}<2.
\]
With one long path, open it at an internal vertex; the complement is an
attached $\Theta(1,2,2)$ and pays for the deleted tree.  With no long path the
core is $K_4$.  We include the attached-$K_4$ argument rather than infer it
from the bare spectrum.  For any graph $H$ obtained by attaching rooted trees
at the four vertices of $K_4$, put
$Z_F(t)=\sum_jm_j(F)t^{|V(F)|-2j}>0$, where $m_j(F)$ counts the $j$-edge
matchings of $F$.  Grouping Sachs' expansion by its vertex-disjoint cycle
collection gives
\[
 \Psi_H(t):=i^{-|V(H)|}\det(itI-A(H))
 =\sum_{\mathcal C}\prod_{C\in\mathcal C}(-2i^{-|C|})
 Z_{H-V(\mathcal C)}(t).
\]
Every cycle lies in the $K_4$, so no two are vertex-disjoint.  The empty term
and the $4$-cycle terms are real, whereas each of the four triangles has
factor $-2i$.  Therefore
\[
 \operatorname{Im}\Psi_H(t)
 =-2\sum_{T\subseteq K_4,\ T\cong K_3}Z_{H-V(T)}(t)<0
 \qquad(t>0).
\]
Moreover $\Psi_H(t)=\prod_j(t+i\lambda_j)$, so
$\sum_j\arctan(\lambda_j/t)$ is its continuous argument tending to zero at
infinity.  Since the curve remains in the open lower half-plane, that
argument is the principal one and is negative.  The signed Coulson identity
\eqref{eq:coulson} yields $D(H)>0$.  As $H$ is tricyclic,
\eqref{eq:energy} gives $\splus(H)>|V(H)|+2$, in particular the required
claim.  This is the attached-$K_4$ lower-half-plane Sachs packet from
\cite{PackingTwo}; bare $K_4$ has spectrum $\{3,-1,-1,-1\}$ and $\splus=9$.
Thus all eight classes are covered.
\end{proof}

\begin{proof}[Proof of Theorem~\ref{thm:main}]
Apply Proposition~\ref{prop:blocks}.  Block ranks $1+1+1$ are covered by
Proposition~\ref{prop:111}, and ranks $2+1$ by Proposition~\ref{prop:21}.
For one rank-three block, the suppressed kernel is one of the four listed in
Proposition~\ref{prop:blocks}; apply Proposition~\ref{prop:fourpath} or
Lemmas~\ref{lem:dc4}, \ref{lem:dtri}, and \ref{lem:k4}.  The DNN additivity
and each structural deletion explicitly include every bridge and rooted-tree
attachment.  The alternatives are exhaustive, proving $\splus(G)\geq n$.
\end{proof}

\appendix
\section{Complete doubled-$C_4$ Gram ledger}\label{app:dc4-ledger}

This appendix makes Lemma~\ref{lem:dc4} reconstructible without consulting
the proof note.  The path order is $(a,A,b,c,C,d)$ and a parity word uses zero
for even and one for odd.  In Table~\ref{tab:dc4-complete}, $\ell$ is the
canonical physical length vector (replace zero by two and one by one), and
$k=(k_1,k_2,k_3)$ means that the four branch vectors have arguments
$\pi(0,k_1,k_2,k_3)/6$.  Thus
\[
 R_{ij}=\cos\frac{\pi(k_i-k_j)}6.
\]
For a path joining $i$ to $j$, let $q\in\{0,\ldots,6\}$ be the principal
integer satisfying
\[
 \arccos((-1)^\ell R_{ij})=\frac{\pi q}{6}.
\]
The six entries in the $q$ column follow the path order, and the certified
excess is therefore reconstructed, with no numerical optimization, as
\begin{equation}\label{eq:dc4-ledger-cost}
 E=\sum_{r=1}^6\ell_r\tan^2\frac{\pi q_r}{12\ell_r}.
\end{equation}
Every displayed Gram matrix is planar and hence positive semidefinite with
unit diagonal.  The rational number in the last column is a strict upper
bound unless equality is printed.

\begin{landscape}
\begin{longtable}{ccccc}
\caption{The complete $28$-row doubled-$C_4$ DNN Gram ledger.}
\label{tab:dc4-complete}\\
\toprule
parity row&canonical $\ell$&$(k_1,k_2,k_3)$&$(q_a,q_A,q_b,q_c,q_C,q_d)$&exact bound for $E$\\
\midrule
\endfirsthead
\toprule
parity row&canonical $\ell$&$(k_1,k_2,k_3)$&$(q_a,q_A,q_b,q_c,q_C,q_d)$&exact bound for $E$\\
\midrule
\endhead
000000&222222&$(0,0,0)$&$(0,0,0,0,0,0)$&$E=0$\\
001001&221221&$(0,6,6)$&$(0,0,0,0,0,0)$&$E=0$\\
000001&222221&$(11,8,7)$&$(1,1,3,1,1,1)$&$E<3/5$\\
001000&221222&$(1,8,9)$&$(1,1,1,1,1,3)$&$E<3/5$\\
000010&222212&$(1,3,11)$&$(1,1,2,4,2,1)$&$E<5/4$\\
000100&222122&$(1,3,11)$&$(1,1,2,2,4,1)$&$E<5/4$\\
010000&212222&$(8,9,10)$&$(4,2,1,1,1,2)$&$E<5/4$\\
100000&122222&$(8,9,10)$&$(2,4,1,1,1,2)$&$E<5/4$\\
000011&222211&$(1,2,6)$&$(1,1,1,4,2,0)$&$E<9/8$\\
000101&222121&$(1,2,6)$&$(1,1,1,2,4,0)$&$E<9/8$\\
001010&221212&$(11,5,1)$&$(1,1,0,4,2,1)$&$E<9/8$\\
001100&221122&$(11,5,1)$&$(1,1,0,2,4,1)$&$E<9/8$\\
010001&212221&$(8,7,6)$&$(4,2,1,1,1,0)$&$E<9/8$\\
100001&122221&$(8,7,6)$&$(2,4,1,1,1,0)$&$E<9/8$\\
011000&211222&$(8,2,1)$&$(4,2,0,1,1,1)$&$E<9/8$\\
101000&121222&$(8,2,1)$&$(2,4,0,1,1,1)$&$E<9/8$\\
001011&221211&$(2,9,5)$&$(2,2,1,4,2,1)$&$E<3/2$\\
001101&221121&$(2,9,5)$&$(2,2,1,2,4,1)$&$E<3/2$\\
011001&211221&$(4,9,7)$&$(4,2,1,2,2,1)$&$E<3/2$\\
101001&121221&$(4,9,7)$&$(2,4,1,2,2,1)$&$E<3/2$\\
010010&212212&$(8,8,0)$&$(4,2,0,4,2,0)$&$E=2$\\
010100&212122&$(8,8,0)$&$(4,2,0,2,4,0)$&$E=2$\\
100010&122212&$(8,8,0)$&$(2,4,0,4,2,0)$&$E=2$\\
100100&122122&$(8,8,0)$&$(2,4,0,2,4,0)$&$E=2$\\
011011&211211&$(4,10,6)$&$(4,2,0,4,2,0)$&$E=2$\\
011101&211121&$(4,10,6)$&$(4,2,0,2,4,0)$&$E=2$\\
101011&121211&$(4,10,6)$&$(2,4,0,4,2,0)$&$E=2$\\
101101&121121&$(4,10,6)$&$(2,4,0,2,4,0)$&$E=2$\\
\bottomrule
\end{longtable}
\end{landscape}

For a paper-only check of the last column, substitute the integer angles in
\eqref{eq:dc4-ledger-cost}, reduce by the half-angle identity
$\tan^2(x/2)=(1-\cos x)/(1+\cos x)$, and use
\[
 \cos\frac\pi6=\frac{\sqrt3}{2},\quad
 \cos\frac\pi4=\frac{\sqrt2}{2},\quad
 \cos\frac\pi3=\frac12.
\]
Clearing positive denominators and squaring positive sides proves the stated
rational comparisons.  In each of the eight equality rows the six terms are,
up to exchanging members within a doubled pair,
\[
 (2/3,1/3,0,2/3,1/3,0),
\]
so equality is exact.  The ledger contains $2+2+4+4+4+4+8=28$ rows, namely
all simple physical rows in switching classes $000$ through $110$.  The eight
remaining simple rows are
\[
010011,\ 010101,\ 011010,\ 011100,\quad
100011,\ 100101,\ 101010,\ 101100;
\]
these are precisely class $111$ and are closed by the long-path certificates
and structural deletion in the proof of Lemma~\ref{lem:dc4}.  Thus the displayed
$28+8=36$ count is the entire simple physical parity census.

\section{Complete doubled-triangle rational ledger}\label{app:dtri-ledger}

We similarly expand the twelve-orbit table used in Lemma~\ref{lem:dtri}.  For
a doubled side of type $S\in\{EE,EO,OO\}$ and rational parameter $t$, put
\begin{align*}
r_{EE}(t)=r_{EO}(t)&=\frac{1-6t^2+t^4}{(1+t^2)^2},
&e_{EE}(t)&=4t^2,\\
&&e_{EO}(t)&=\left(\frac{1-t^2}{2t}\right)^2+2t^2,\\
q(t)&=\frac{3t-t^3}{1-3t^2},
&r_{OO}(t)&=-\frac{1-q(t)^2}{1+q(t)^2},
&e_{OO}(t)&=q(t)^2+3t^2.
\end{align*}
For the connector define
\[
r_E(t)=\frac{1-6t^2+t^4}{(1+t^2)^2},\quad e_E(t)=2t^2,
\qquad r_O(t)=\frac{t^2-1}{t^2+1},\quad e_O(t)=t^2.
\]
Given $(t_1,t_2,t_c)$, the table uses
\[
R=\begin{pmatrix}1&r_1&r_2\\r_1&1&r_c\\r_2&r_c&1\end{pmatrix},
\qquad E=e_{S_1}(t_1)+e_{S_2}(t_2)+e_P(t_c),
\]
where $P$ is the connector parity.  These multiple-angle formulas reconstruct
every parameterized row using rational arithmetic only.  The three boundary
rows use their displayed correlations directly.

\begin{landscape}
\begin{longtable}{ccccccc}
\caption{Complete doubled-triangle rational certificate table.  The orbit
sizes total $32$ labelled physical rows.}\label{tab:dtri-complete}\\
\toprule
types&$c$&size&$(t_1,t_2,t_c)$&$(r_1,r_2,r_c)$&$E$&$\det R$\\
\midrule
\endfirsthead
\toprule
types&$c$&size&$(t_1,t_2,t_c)$&$(r_1,r_2,r_c)$&$E$&$\det R$\\
\midrule
\endhead
$EE,EE$&$E$&1&$(0,0,0)$&$(1,1,1)$&$0$&$0$\\
$EE,EE$&$O$&1&$(1/3,1/3,1/2)$&$(7/25,7/25,-3/5)$&$41/36$&$1216/3125$\\
$EE,EO$&$E$&4&$(0,1/2,1/2)$&$(1,-7/25,-7/25)$&$25/16$&$0$\\
$EE,EO$&$O$&4&$(1/4,1/2,1/3)$&$(161/289,-7/25,-4/5)$&$205/144$&$11527191/52200625$\\
$EE,OO$&$E$&2&$(1/4,1/4,1/3)$&$(161/289,-495/4913,7/25)$&$2246/1521$&$8593933244/15085980625$\\
$EE,OO$&$O$&2&boundary&$(1,-1,-1)$&$0$&$0$\\
$EO,EO$&$E$&4&boundary&$(-1/2,-1/2,1)$&$2$&$0$\\
$EO,EO$&$O$&4&class $111$&---&structural&---\\
$EO,OO$&$E$&4&$(1/2,1/5,1/5)$&$(-7/25,-828/2197,119/169)$&$83009/48400$&$1304316959/3016755625$\\
$EO,OO$&$O$&4&$(1/2,1/5,1/2)$&$(-7/25,-828/2197,-3/5)$&$91237/48400$&$883705599/3016755625$\\
$OO,OO$&$E$&1&boundary&$(-1,-1,1)$&$0$&$0$\\
$OO,OO$&$O$&1&$(1/5,1/5,1/2)$&$(-828/2197,-828/2197,-3/5)$&$16881/12100$&$22382224/120670225$\\
\bottomrule
\end{longtable}
\end{landscape}

For each nonstructural row, $|r_1|,|r_2|,|r_c|\leq1$ and the displayed
determinant is nonnegative; hence all principal minors of $R$ are nonnegative
and $R\succeq0$.  Direct substitution gives the displayed excess fraction.
The only less immediate threshold checks are
\[
2246<2(1521),\quad 83009<2(48400),\quad
91237<2(48400),\quad 16881<2(12100),
\]
and every other DNN row visibly has $E\leq2$.  The size column sums to
$16$ rows for each connector parity, or $32$ in all; the structural orbit has
size four, leaving $28$ labelled DNN rows.  This supplies all correlations,
excesses, determinants, and multiplicities needed to reconstruct the finite
part of the doubled-triangle proof from the paper alone.

\section{Reproducibility, dependencies, and scope}

The detailed exact tables summarized above are repository proof objects, not
empirical graph censuses.  They prove arbitrary path lengths by the monotonic
path lemma and arbitrary rooted trees by one-vertex DNN additivity or explicit
territory ownership.  The source notes are:
\begin{verbatim}
positive-square-energy/tricyclic-general/four-path-dnn-theorem.md
positive-square-energy/tricyclic-general/theta-cycle-dnn-reduction.md
positive-square-energy/tricyclic-general/theta-cycle-completion.md
positive-square-energy/tricyclic-general/doubled-c4-switching-sieve.md
positive-square-energy/tricyclic-general/k4-seven-switching-classes.md
positive-square-energy/tricyclic-general/k4-all-odd-dnn-cover.md
positive-square-energy/tricyclic-general/doubled-triangle-dnn-cover.md
\end{verbatim}
The finite rational audit checks twelve doubled-triangle type/parity orbits
(32 labelled physical rows), partitioned accurately as 28 DNN rows and four
structural class-$111$ rows; it also checks the two exact rational tangent
certificates used when a class-$111$ parallel path is noncanonical.  It checks
the doubled-$4$-cycle partition into 28 DNN rows and eight structural
class-$111$ rows, all 56 physical $K_4$ rows outside the all-odd switching
class, and explicitly enumerates all 64 all-odd long/unit subsets.  The latter
split into 42 DNN-audited subsets with at least three long paths, 12 adjacent
and three opposite DNN-audited two-long subsets, six one-long structural
deletions, and one no-long attached-$K_4$ packet.  Thus 57 subsets are audited
by the displayed DNN bounds and seven have structural dispositions.  The audit
also checks nine four-path representative symbolic records, including the
one-unit $e=0$ endpoint case, and the rational tangent bounds occurring in
those records.  These records do not constitute a machine proof of the
analytic tangent lemma or its all-length coverage.  The verifier rejects
thirteen deliberately corrupted certificate payloads.
For the 56-row audit it reconstructs every switching invariant and physical
sign transport while retaining the physical path parity.  Run the audit both
normally and with Python assertions disabled:
\begin{verbatim}
python3 research/tricyclic-finite-rational-certificates-verifier.py
python3 -O research/tricyclic-finite-rational-certificates-verifier.py
\end{verbatim}
Both runs must report the certificate digest
\begin{verbatim}
a34ed2d3898c1a244e861ce11ffb51d84d65eb4409066f0745829de5a8ca58b8
\end{verbatim}
This canonical payload digest is embedded in the verifier and mirrored in the
README manifest.  Both runs must also report thirteen rejected hostile
mutations.  More generally, the verifier does not
prove DNN duality, path elimination or monotonicity, induced
superadditivity, the kernel classification, the territory repairs, or the
attached-$K_4$ Sachs argument.  It is a fail-closed exact audit of only the
finite records just enumerated.
Build this manuscript from the repository root with
\begin{verbatim}
bash scripts/build-paper.sh all-tricyclic-graphs
\end{verbatim}

\paragraph{Scoped nonclaim.}
The theorem concerns exactly finite simple connected graphs with $m=n+2$.
Some branches above are strict, but DNN certificates may meet their auxiliary
boundary.  We therefore claim only $\splus(G)\geq n$ globally.  We do not claim
that all DNN equality cases have been spectrally classified, nor do we invoke
edge monotonicity or the false inequality $\splus(G)\geq\operatorname{MaxCut}(G)$.

\section*{AI disclosure}
This manuscript and its proof-object synthesis were generated by an AI coding
and mathematical reasoning system from the cited repository materials.  No
human authorship, review, or verification is claimed.  The result should be
assessed from the displayed proof and its exact local dependencies.  No
external publication was performed as part of this work.

\begin{thebibliography}{99}
\bibitem{BicyclicCactus}
Companion manuscript, \emph{Positive Square Energy of Every Bicyclic Cactus},
2026; \path{all-bicyclic-cacti/paper.tex}.

\bibitem{NonunicyclicCactus}
Companion manuscript, \emph{Positive Square Energy of Every Nonunicyclic
Cactus}, 2026; \path{all-nonunicyclic-cacti/paper.tex}.

\bibitem{DoubledC4}
\emph{The doubled-$C_4$ kernel: exact switching sieve}, proof note, 2026;
\path{positive-square-energy/tricyclic-general/doubled-c4-switching-sieve.md}.

\bibitem{DoubledTriangle}
\emph{The doubled-triangle kernel: complete physical parity cover}, proof note,
2026; \path{positive-square-energy/tricyclic-general/doubled-triangle-dnn-cover.md}.

\bibitem{K4Cover}
\emph{The all-odd $K_4$-subdivision class}, proof note, 2026;
\path{positive-square-energy/tricyclic-general/k4-all-odd-dnn-cover.md}.

\bibitem{K4Seven}
\emph{The seven non-all-odd switching classes of subdivided $K_4$}, proof
note, 2026;
\path{positive-square-energy/tricyclic-general/k4-seven-switching-classes.md}.

\bibitem{PackingTwo}
Companion manuscript, \emph{Square-energy sign under cycle packing number at
most two}, 2026; \path{packing-two-square-energy/paper.tex}.
\end{thebibliography}

\end{document}
